%% kapa: exam=1 round=32 year=2021 period=1 form=A subject=civil_law subject_ko=민법 seq=6 number=6
%% source: Q-Net 감정평가사 기출문제 (cst00309) artlSeq=5208397 file=extracted/1008/5208397/A형/2021년도 제32회 감정평가사 1차 1교시 A형.pdf
%% conversion: kordoc 4.24.51 -> tools/gen1.py (auto); review=proofread-sonnet; flags=none
\begin{question}{6}
법인의 불법행위책임에 관한 설명으로 옳은 것은? (다툼이 있으면 판례에 따름)
\begin{choices}
  \item 외형상 직무행위로 인정되는 대표자의 권한 남용행위에 대해서도 법인의 불법행위책임이 인정될 수 있다.
  \item 등기된 대표자의 행위로 인하여 타인에게 손해를 가한 경우에만 법인의 불법행위책임이 성립할 수 있다.
  \item 대표자의 행위가 직무에 관한 행위에 해당하지 않음을 피해자 자신이 중대한 과실로 알지 못한 경우, 법인의 불법행위책임이 인정된다.
  \item 대표권 없는 이사가 그 직무와 관련하여 타인에게 손해를 가한 경우, 법인의 불법행위책임이 성립한다.
  \item 법인의 불법행위책임이 성립하는 경우 그 대표기관은 손해배상책임이 없다.
\end{choices}
\end{question}
\answer{1}
